% !TEX root = ../../report.tex

\section{Анализ существующих решений}

Для определения функциональных требований к проектируемой системе проведён сравнительный анализ трёх распространённых маркетплейсов, доступных российским пользователям.

\begin{itemize}
    \item \textbf{Ozon} --- маркетплейс с интерфейсом для продавцов, покрывающим управление товарами, заказами, аналитикой и финансами. Предоставляет расширенную аналитику в личном кабинете продавца, часть функций доступна по подписке. Поддерживает раздельные рейтинги товара и продавца~\cite{ozon};
    \item \textbf{Wildberries} --- маркетплейс с личным кабинетом продавца, охватывающим товары, заказы, поставки и продвижение. Базовая аналитика встроена в платформу бесплатно; для расширенной требуются сторонние сервисы. Поддерживаются раздельные рейтинги товара и продавца~\cite{wildberries};
    \item \textbf{AliExpress} --- трансграничная торговая площадка, ориентированная на международных продавцов. Базовая аналитика доступна в кабинете продавца. Поддерживаются рейтинги товара и продавца~\cite{aliexpress}.
\end{itemize}

Критерии сравнения соответствуют характеристикам, которые используются в проектируемой системе:

\begin{enumerate}
    \item открытость модели данных;
    \item история изменения цен товара;
    \item бесплатная встроенная аналитика продаж;
    \item раздельные рейтинги товара и продавца;
    \item выделенная роль аналитика только для чтения;
    \item логистика как часть предметной области.
\end{enumerate}

Сравнение платформ по сформированным критериям приведено в таблице~\ref{tbl:comparison}.

\begin{table}[h!]
    \caption{Сравнение существующих решений}
    \label{tbl:comparison}
    \centering
    \small
    \begin{tabular}{|l|l|l|l|l|}
        \hline
        \textbf{Критерий} & \textbf{Ozon} & \textbf{Wildberries} & \textbf{AliExpress} & \makecell[l]{\textbf{Проектируемая}\\\textbf{система}} \\ \hline
        \makecell[l]{Открытость\\модели данных} & Нет & Нет & Нет & Да \\ \hline
        \makecell[l]{История\\цен товара} & \makecell[l]{Не\\раскрывается} & \makecell[l]{Не\\раскрывается} & \makecell[l]{Не\\раскрывается} & Да \\ \hline
        \makecell[l]{Бесплатная встроенная\\аналитика продаж} & Частично & Да & Да & Да \\ \hline
        \makecell[l]{Рейтинг товара\\и продавца} & Да & Да & Да & Да \\ \hline
        \makecell[l]{Выделенная роль\\аналитика\\только для чтения} & \makecell[l]{Не\\указана} & \makecell[l]{Не\\указана} & \makecell[l]{Не\\указана} & Да \\ \hline
        \makecell[l]{Логистика\\в предметной\\области} & Да & Да & Да & Нет \\ \hline
    \end{tabular}
\end{table}

Рассмотренные платформы являются закрытыми коммерческими системами.

В проектируемой системе выделяются открытая модель данных, хранение истории изменения цен и отдельная роль аналитика. Задача логистики не рассматривается.
